% ============ 第2章 §2.5 Hilbert 空间 (原书 pp.54-60) ============

\section{Hilbert 空间}\label{sec:2.5}
\index{0symbolAperp@$A^{\perp}$}

\index{0symbolAoplus@$A_{1} \oplus A_{2}$}

\index{0symbolAperp2@$A_{1} \perp A_{2}$}

\index{C@Cauchy--Schwarz 不等式}

\index{Z@准 Hilbert 空间}

\index{Z@正交直和}

\index{Y@有界自伴算子}

\index{Y@酉线性算子}


















调和分析学中, 算子不是 $L^{1}$-有界而是 $L^{2}$-有界的情形很多. 又, 利用 $L^{2}$-有界性导出 $L^{p}$-有界性的状况很多, 本书也会遭遇这种场面. 因此 $L^{2}$ 空间重要, 理解 $L^{2}$ 空间所具有的 Hilbert 空间结构也是重要的. 这里总结 Hilbert 空间. 附上证明列举 Hilbert 空间的基础知识.

\subsection*{准 Hilbert 空间}

一度忘掉完备性, 总结可代数处理的部分.

\begin{definition}\label{def:2.21}
设 $H$ 为复线性空间. $\langle \cdot, \cdot \rangle_{\mathcal{H}}: H \times H \to \mathbb{C}$ 是 Hermitian 内积, 是指对所有 $u, u_{1}, u_{2}, v, v_{1}, v_{2} \in H$ 与 $\alpha \in \mathbb{C}$ 满足以下条件.

(1) 若 $u \neq 0$, 则 $\langle u, u \rangle_{\mathcal{H}} > 0$. 【内积的正值性】

(2) $\langle \alpha u, v \rangle_{\mathcal{H}} = \langle u, \overline{\alpha}v \rangle_{\mathcal{H}} = \alpha\langle u, v \rangle_{\mathcal{H}}$.

(3) $\langle u, v \rangle_{\mathcal{H}} = \overline{\langle v, u \rangle_{\mathcal{H}}}$.

(4) $\langle u_{1} + u_{2}, v \rangle_{\mathcal{H}} = \langle u_{1}, v \rangle_{\mathcal{H}} + \langle u_{2}, v \rangle_{\mathcal{H}}$, $\langle u, v_{1} + v_{2} \rangle_{\mathcal{H}} = \langle u, v_{1} \rangle_{\mathcal{H}} + \langle u, v_{2} \rangle_{\mathcal{H}}$.

配备 Hermitian 内积的线性空间称为 (复) 准 Hilbert 空间. 又, 在上述定义中把 $\mathbb{C}$ 换成 $\mathbb{R}$ 而得到的运算 $\langle \cdot, \cdot \rangle_{\mathcal{H}}: \mathcal{H} \times \mathcal{H} \to \mathbb{R}$ 所配备的线性空间称为 (实) 准 Hilbert 空间. 又, 令 $\|x\| = \sqrt{\langle x, x \rangle_{\mathcal{H}}}$, $x \in \mathcal{H}$. (2), (3), (4) 称为双线性性.
\end{definition}

举一个用准 Hilbert 空间中的运算计算的例子.

\begin{example}\label{exa:2.5}
设 $H$ 为准 Hilbert 空间, $k \in \mathbb{C}$, $x, y \in H$. 此时,
\begin{align*}
\|x + ky\|^{2} &= \langle x + ky, x + ky \rangle_{\mathcal{H}} \\
&= \langle x, x \rangle_{\mathcal{H}} + \langle x, ky \rangle_{\mathcal{H}} + \langle ky, x \rangle_{\mathcal{H}} + \langle ky, ky \rangle_{\mathcal{H}} \\
&= \|x\|^{2} + \overline{k}\langle x, y \rangle_{\mathcal{H}} + k\overline{\langle x, y \rangle_{\mathcal{H}}} + |k|^{2}\|y\|^{2} \\
&= \|x\|^{2} + 2\mathrm{Re}(\overline{k}\langle x, y \rangle_{\mathcal{H}}) + |k|^{2}\|y\|^{2}
\end{align*}
成立.
\end{example}

准 Hilbert 空间的例子有下面这个. $L^{2}(\mathbb{R}^{n})$ 的内积取为
\[
\langle f, g \rangle_{L^{2}} = \int_{\mathbb{R}^{n}}f(x)\overline{g(x)}\dd x
\]
时, $L^{2}(\mathbb{R}^{n})$ 是准 Hilbert 空间. 不过, 如上所述, $L^{2}(\mathbb{R}^{n})$ 配备有完备距离, 所以由后述定义可知 $L^{2}(\mathbb{R}^{n})$ 是 Hilbert 空间. 本书不考虑是准 Hilbert 空间但不是 Hilbert 空间的例子, 所以只举这个例子.

\subsection*{Cauchy--Schwarz 不等式}

证明从 Hilbert 空间的性质可代数地看出的 Cauchy--Schwarz 不等式. 由此可知 $H$ 上可定义距离.

\begin{theorem}[Cauchy--Schwarz 不等式]\label{thm:2.25}
$H$ 为准 Hilbert 空间时, $|\langle x, y \rangle_{H}| \leq \|x\|\|y\|$, $x, y \in H$ 成立.
\end{theorem}

\begin{proof}
不妨设 $x, y \in H \setminus \{0\}$. 令 $\alpha = -\frac{\langle x, y \rangle_{H}}{\|y\|^{2}} \in \mathbb{C}$,
\[
\|x + \alpha y\|^{2} = \|x\|^{2} + 2\mathrm{Re}(\overline{\alpha}\langle x, y \rangle_{\mathcal{H}}) + |\alpha|^{2}\|y\|^{2} = \|x\|^{2} - \frac{|\langle x, y \rangle_{\mathcal{H}}|^{2}}{\|y\|^{2}}
\]
由于这个量非负, 所以显然.
\end{proof}

\subsection*{Hilbert 空间}

给定准 Hilbert 空间后, 由 Cauchy--Schwarz 不等式可知范数被确定, 但如前所述, 关于该范数完备时另有名称. 下面认真表述它.

\begin{definition}\label{def:2.22}
(复) Hilbert 空间是指配备 Hermitian 内积 $\langle \cdot, \cdot \rangle = \langle \cdot, \cdot \rangle_{H}$ 的内积空间, 且关于由此诱导的距离 $\|\cdot\| = \|\cdot\|_{H} = \sqrt{\langle \cdot, \cdot \rangle_{H}}$ 完备者.
\end{definition}

$H$ 的例子有 $L^{2}(\mathbb{R}^{n})$, $\ell^{2}(\mathbb{N})$. 此后设 $H$ 为 Hilbert 空间.

Hilbert 空间是 Banach 空间, 所以有界线性算子等概念自然可以设想, 但与 Hilbert 空间固有配备的内积相关联的算子变得重要. 重要的例子后面给出若干, 这里列举有界线性算子中的酉 (线性) 算子. $H$ 的线性算子 $U$ 是酉 (线性) 算子, 是指 $U$ 为双射且对所有 $x \in H$,
\begin{equation}\tag{2.7}
\|Ux\| = \|x\|
\label{eq:2.7}
\end{equation}
成立. 例如, 对 $\ell^{2}(\mathbb{N})$, 把 $\{a_{n}\}_{n=1}^{\infty} = (a_{1}, a_{2}, \ldots)$ 移到 $(a_{2}, a_{1}, a_{4}, a_{3}, a_{6}, a_{5}, \ldots)$ 的算子是酉算子.

此后, 设 $H$ 为 Hilbert 空间.

\subsection*{中线定理}

由代数计算可知下面的等式成立. Hilbert 空间的几何结构被概括在下面的中线定理中. 不作证明, 但已知满足中线定理的范数可用内积表示.

\begin{theorem}[中线定理]\label{thm:2.26}
对 $x, y \in H$, $\|x + y\|^{2} + \|x - y\|^{2} = 2\|x\|^{2} + 2\|y\|^{2}$ 成立.
\end{theorem}

\begin{proof}
由例~\ref{exa:2.5}, $\|x + y\|^{2} = \|x\|^{2} + \|y\|^{2} + 2\mathrm{Re}(\langle x, y \rangle_{\mathcal{H}})$ 与 $\|x - y\|^{2} = \|x\|^{2} + \|y\|^{2} - 2\mathrm{Re}(\langle x, y \rangle_{\mathcal{H}})$ 显然.
\end{proof}

\subsection*{Hilbert 空间上的有界线性算子}

设 $A$ 为从 $H$ 到 $H$ 的有界线性算子时, $\|A\|$ 已在定义~\ref{def:2.20} 中定义, 但给出由 $H$ 是 Hilbert 空间得到的另一个表示.

\begin{theorem}\label{thm:2.27}
对 $A \in B(\mathcal{H})$, $\|A\|_{\mathcal{H} \to \mathcal{H}} = \sup_{x, y \in \mathcal{H}, \|x\|, \|y\| \leq 1}|\langle Ax, y \rangle_{\mathcal{H}}|$ 成立.
\end{theorem}

\begin{proof}
若 $A = 0$, 则结论显然, 所以不妨设 $A \neq 0$. 由 Cauchy--Schwarz 不等式, 对 $x, y \in H$,
\[
|\langle Ax, y \rangle_{\mathcal{H}}| \leq \|Ax\|\|y\| \leq \|A\|_{\mathcal{H} \to \mathcal{H}}\|x\|\|y\|
\]
所以
\[
\|A\|_{\mathcal{H} \to \mathcal{H}} \geq \sup_{x, y \in \mathcal{H}, \|x\|, \|y\| \leq 1}|\langle Ax, y \rangle_{\mathcal{H}}|
\]
成立. 又,
\[
\sup_{x, y \in \mathcal{H}, \|x\|, \|y\| \leq 1}|\langle Ax, y \rangle_{\mathcal{H}}| = \sup_{x, y \in \mathcal{H} \setminus \{0\}}\frac{|\langle Ax, y \rangle_{\mathcal{H}}|}{\|x\|\|y\|}
\]
所以, 特别是对 $Ax \neq 0$ 的 $x \in H$, 通过令 $y = Ax$ 可部分到达上限,
\[
\frac{\|Ax\|}{\|x\|} = \frac{|\langle Ax, y \rangle_{\mathcal{H}}|}{\|x\|\|y\|}
\]
成立. 因此,
\[
\|A\|_{\mathcal{H} \to \mathcal{H}} \leq \sup_{x, y \in \mathcal{H}, \|x\|, \|y\| \leq 1}|\langle Ax, y \rangle_{\mathcal{H}}|
\]
成立.
\end{proof}

\subsection*{最短距离}

下面的定理保证图形的最短距离可达性. 有限维情形由于可用紧性, 结论显然. 但无限维情形不保证单位闭球的紧性 (参照练习问题~\ref{exer:1.29}), 即单位闭球不是全有界的, 所以证明需要技巧. 作为全有界性的代用物的是中线定理. 下面定理中出现的凸集合是指, 任意取该集合内的 $2$ 点, 以这 $2$ 点为端点的线段包含在该集合中.

\begin{theorem}\label{thm:2.28}
闭 (且) 凸集合 $C$ 与 $H$ 内的点 $x$ 的距离被 $C$ 内的点 $y$ 唯一地实现. 即, 存在唯一的 $y \in C$ 使得 $\|x - y\| = \inf_{z \in C}\|x - z\|$ 成立.
\end{theorem}

\begin{proof}
令 $d = \inf_{z \in C}\|x - z\|$. 对 $j \in \mathbb{N}$, 取 $y_{j} \in C$ 使得 $\|x - y_{j}\| \leq d + 2^{-j}$ 成立. 则由中线定理, 对各 $j, k \in \mathbb{N}$,
\[
\|2x - y_{j} - y_{k}\|^{2} + \|y_{j} - y_{k}\|^{2} = 2\|x - y_{j}\|^{2} + 2\|x - y_{k}\|^{2}
\]
成立. 这里, 由于 $C$ 是凸的, $\frac{y_{j} + y_{k}}{2} \in C$. 因此 $\|2x - y_{j} - y_{k}\| \geq 2d$, 得到
\[
\|y_{j} - y_{k}\|^{2} \leq 2(d + 2^{-j})^{2} + 2(d + 2^{-k})^{2} - 4d^{2}.
\]
令 $j, k \to \infty$, 则 $y_{j} - y_{k} \to 0$, 所以由 $H$ 的完备性, $\{y_{j}\}_{j=1}^{\infty}$ 收敛到某个 $y$. 因此 $y \in C$ 是所求的点.

假若 $y' \in C$ 也给出距离. 由于 $\frac{y + y'}{2} \in C$, $\|x - y'\| = d$,
\[
4d^{2} \leq \|2x - y - y'\|^{2} + \|y - y'\|^{2} = 2\|x - y\|^{2} + 2\|x - y'\|^{2} = 4d^{2}
\]
得到 $\|y - y'\| = 0$, 即 $y = y'$.
\end{proof}

\subsection*{正交补空间}

对 $H$ 的子集 $A \subset H$, 令 $A^{\perp} = \{x \in H : \langle x, z \rangle_{H} = 0, z \in A\}$. 对 $H$ 的子集 $A_{1}$, $A_{2}$, $A_{1} \perp A_{2}$ 是指对所有 $a_{1} \in A_{1}$, $a_{2} \in A_{2}$, $a_{1} \perp a_{2}$, 即 $\langle a_{1}, a_{2} \rangle_{H} = 0$ 成立. 进而, 当 $A_{1}$, $A_{2}$ 为线性空间时, 把和空间
\[
A_{1} + A_{2} = \{a_{1} + a_{2} \in \mathcal{H} : a_{1} \in A_{1}, a_{2} \in A_{2}\}
\]
表示为 $A_{1} \oplus A_{2}$. 这个 $A_{1} \oplus A_{2}$ 称为正交直和.

\begin{theorem}\label{thm:2.29}
$Z$ 是 $H$ 的闭线性子空间时, $H = Z \oplus Z^{\perp}$ 成立.
\end{theorem}

\begin{proof}
由内积的正值性, $Z \cap Z^{\perp} = \{0\}$, $H \supset Z + Z^{\perp}$, $Z \perp Z^{\perp}$, 即 $H \supset Z \oplus Z^{\perp}$ 显然. 因此必须证明 $H \subset Z \oplus Z^{\perp}$. 设 $x \in H$. 由定理~\ref{thm:2.28}, 存在 $z_{0} \in Z$ 使得 $z_{0}$ 给出 $x$ 与 $Z$ 之间的距离, 所以对所有 $z \in Z$, $\|x - z_{0}\| \leq \|x - z\|$ 成立. 通过把 $z$ 换成 $z_{0} + z$, 此条件与对 $z \in Z$ 有 $\|x - z_{0}\| \leq \|x - z_{0} - z\|$ 成立等价. 因此由例~\ref{exa:2.5},
\[
\|z\|^{2} \geq 2\mathrm{Re}(\langle x - z_{0}, z \rangle_{\mathcal{H}})
\]
成立. 通过把 $z$ 换成 $\kappa z$, $\kappa > 0$, 可知
\[
\kappa\|z\|^{2} \geq 2\mathrm{Re}(\langle x - z_{0}, z \rangle_{\mathcal{H}})
\]
成立. 令 $\kappa \downarrow 0$, 得 $0 \geq 2\mathrm{Re}(\langle x - z_{0}, z \rangle_{\mathcal{H}})$ 成立. 由于 $z$ 也可以换成 $-z$, $0 \leq 2\mathrm{Re}(\langle x - z_{0}, z \rangle_{\mathcal{H}})$ 也成立. 因此 $2\mathrm{Re}(\langle x - z_{0}, z \rangle_{\mathcal{H}}) = 0$. 由于 $z$ 也可以换成 $iz$, $2\mathrm{Im}(\langle x - z_{0}, z \rangle_{\mathcal{H}}) = 0$ 成立. 由这两个等式, 得 $\langle x - z_{0}, z \rangle_{\mathcal{H}} = 0$. 综上, 得 $x - z_{0} \perp z$. 因此 $x = z_{0} + (x - z_{0}) \in Z \oplus Z^{\perp}$, 结论得证.
\end{proof}

由代数计算可知, $\perp$ 取两次会回到原处.

\begin{corollary}\label{cor:2.1}
$Z$ 是 $H$ 的闭线性子空间时, $(Z^{\perp})^{\perp} = Z$ 成立.
\end{corollary}

\begin{proof}
$\mathcal{H} = Z \oplus Z^{\perp} = (Z^{\perp})^{\perp} \oplus Z^{\perp}$ 成立. 由 $Z^{\perp}$, $(Z^{\perp})^{\perp}$ 的定义, $(Z^{\perp})^{\perp} \supset Z$ 显然. 于是, 为证明 $(Z^{\perp})^{\perp} \subset Z$, 任取 $z \in (Z^{\perp})^{\perp}$. 则 $z = z_{0} + z_{1}$, $z_{0} \in Z$, $z_{1} \in Z^{\perp}$ 的分解是可能的. 因为 $z_{0}$, $z$ 均与 $z_{1}$ 正交, 在 $\langle z_{1}, z_{1} \rangle_{\mathcal{H}}$ 中将右侧的 $z_{1}$ 代入 $z_{1} = z - z_{0}$,
\[
\langle z_{1}, z_{1} \rangle_{\mathcal{H}} = \langle z_{1}, z \rangle_{\mathcal{H}} - \langle z_{1}, z_{0} \rangle_{\mathcal{H}} = 0
\]
所以得到 $z_{1} = 0$. 即 $z = z_{0} \in Z$.
\end{proof}

把系~\ref{cor:2.1} 稍微放宽, 以便使用.

\begin{corollary}\label{cor:2.2}
$Z$ 是不一定是闭子空间的 $H$ 的线性子空间时, $(Z^{\perp})^{\perp} = \overline{Z}$ 成立.
\end{corollary}

\begin{proof}
$\overline{Z}$ 是闭线性空间, 由系~\ref{cor:2.1}, $\overline{Z} = (\overline{Z}^{\perp})^{\perp} = (Z^{\perp})^{\perp}$, 即得.
\end{proof}

\subsection*{Riesz 表示定理}

与许多数学领域一样, 调和分析学中利用对偶性也是基本的. 已经说明的 $L^{p}$ 与 $L^{p'}$ 的对偶性即其一例. 下面证明相当于 Hilbert 空间中的自对偶性的 Riesz 表示定理. 回忆: 对线性空间 $U$, $V$ 与线性映射 $f: U \to V$, $f$ 的核 $\ker(f)$ 定义为 $\ker(f) = \{x \in U : f(x) = 0\}$, $f$ 的像 $\operatorname{im}(f)$ 定义为 $\operatorname{im}(f) = \{f(x) \in V : x \in U\}$. 证明中出现的同构定理是指断言 $\operatorname{im}(f) \simeq U/\ker(f)$ 成立的定理.

\begin{theorem}[Riesz 表示定理]\label{thm:2.30}
从 $H$ 到 $\mathbb{C}$ 的连续线性泛函可借助唯一的 $y \in H$ 由内积表示为 $\langle \cdot, y \rangle_{H}$ 的形式. 即, $H^{*}$ 与 $H$ 被同一化.
\end{theorem}

\begin{proof}
设 $\ell$ 为从 $\mathcal{H}$ 到 $\mathbb{C}$ 的连续线性泛函. $\ker(\ell) = \{x \in \mathcal{H} : \ell(x) = 0\} = \ell^{-1}(\{0\})$ 是 $\mathcal{H}$ 的闭子空间, 所以 $\mathcal{H} = \ker(\ell) \oplus \ker(\ell)^{\perp}$ 成立. 由同构定理, $\dim(\ker(\ell)^{\perp}) = \dim(\mathcal{H}/\ker(\ell)) = \dim(\operatorname{im}(\ell)) \leq 1$. 若 $\dim(\ker(\ell)^{\perp}) = 0$, 则 $\ker(\ell) = \mathcal{H}$, 所以 $\ell = 0$, $y = 0$.

若 $\dim(\ker(\ell)^{\perp}) = 1$, 取 $\ker(\ell)^{\perp}$ 的基底 $\{y_{0}\}$ 使 $\|y_{0}\| = 1$ 成立. 任意元 $x \in \mathcal{H}$ 满足 $x - \langle x, y_{0} \rangle_{\mathcal{H}}y_{0} \in (\ker(\ell)^{\perp})^{\perp} = \ker(\ell)$, 所以 $\ell(x) - \langle x, y_{0} \rangle_{\mathcal{H}}\ell(y_{0}) = 0$. 因此, 令 $y = \overline{\ell(y_{0})}y_{0}$, 得 $\ell(x) = \langle x, y \rangle_{\mathcal{H}}$.

最后, 与 $\dim(\ker(\ell)^{\perp})$ 的值无关, $y \in H$ 的唯一性由内积的正值性显然.
\end{proof}

\subsection*{伴随算子}

这里给出 Riesz 表示定理的应用. 伴随算子与对偶性相关联, 在调和分析学中重要.

\begin{theorem}\label{thm:2.31}
对从 $H$ 到 $H$ 的连续线性算子 $A$, 可以说以下事实.

(1) 称为 $A$ 的伴随算子的从 $H$ 到 $H$ 的连续线性算子 $A^{*}$ 唯一确定, 且 $\langle Ax, y \rangle_{H} = \langle x, A^{*}y \rangle_{H}$ 成立.

(2) $\|A^{*}A\|_{H \to H} = (\|A\|_{H \to H})^{2}$ 成立.

(3) 关于 $\ker(A^{*}) = \{x \in \mathcal{H} : A^{*}x = 0\}$ 与 $\operatorname{im}(A) = \{Ax \in \mathcal{H} : x \in \mathcal{H}\}$, 有关系 $\ker(A^{*}) = \operatorname{im}(A)^{\perp}$, $\ker(A^{*})^{\perp} = \overline{\operatorname{im}(A)}$ 成立.
\end{theorem}

\begin{proof}
(1) 固定 $x \in \mathcal{H}$. 线性映射 $\ell(y) = \langle Ay, x \rangle_{\mathcal{H}}$, $y \in \mathcal{H}$ 是从 $\mathcal{H}$ 到 $\mathbb{C}$ 的连续线性泛函, 所以由 Riesz 表示定理, 存在唯一的 $x' \in \mathcal{H}$ 满足 $\ell(y) = \langle y, x' \rangle_{\mathcal{H}}$. 由内积的性质 $\langle u, v \rangle_{\mathcal{H}} = \overline{\langle v, u \rangle_{\mathcal{H}}}$, $u, v \in \mathcal{H}$, 可知 $x \mapsto x' = A^{*}x$ 满足条件. 最后, 唯一性由内积的正值性显然.

(2) 由定理~\ref{thm:2.27} 与 Cauchy--Schwarz 不等式,
\begin{align*}
\|A^{*}A\|_{\mathcal{H} \to \mathcal{H}} &= \sup_{x, y \in \mathcal{H}, \|x\| \leq 1, \|y\| \leq 1}|\langle A^{*}Ax, y \rangle_{\mathcal{H}}| \\
&= \sup_{x, y \in \mathcal{H}, \|x\| \leq 1, \|y\| \leq 1}|\langle Ax, Ay \rangle_{\mathcal{H}}| \\
&= (\|A\|_{\mathcal{H} \to \mathcal{H}})^{2}.
\end{align*}

(3) 前者由伴随的定义显然. 后者用系~\ref{cor:2.2}, 对前者的关系取正交补空间即可.
\end{proof}

举调和分析中常出现的伴随算子的例子. 对函数 $g$, $\overline{g}$ 是复共轭. 另外, 这个算子的有界性由 Young 不等式得到.

\begin{example}\label{exa:2.6}
$g \in L^{1}(\mathbb{R})$ 时, 有界线性算子 $L^{2}(\mathbb{R}) \ni f \mapsto f * g \in L^{2}(\mathbb{R})$ 的伴随是 $L^{2}(\mathbb{R}) \ni f \mapsto f * \overline{g}(-\cdot) \in L^{2}(\mathbb{R})$. 参照练习问题~\ref{exer:4.4}.
\end{example}

\subsection*{Hilbert 空间上的有界自伴算子}

最后, 总结 Hilbert 空间上的有界自伴算子. 满足 $A = A^{*}$ 的算子 $A \in B(\mathcal{H})$ 称为有界自伴算子, 此类算子有称为谱分解的重要分解理论. 本书不需要谱分解, 但需要有界自伴算子的若干公式, 总结如下.

\begin{theorem}\label{thm:2.32}
设 $H$ 为 Hilbert 空间. 设 $A \in B(\mathcal{H})$ 为自伴算子. 此时, $\|A\|_{\mathcal{H} \to \mathcal{H}} = \sup_{x \in \mathcal{H}, \|x\| \leq 1}|\langle Ax, x \rangle_{\mathcal{H}}|$ 成立.
\end{theorem}

\begin{proof}
对任意范数不超过 $1$ 的 $x, y \in H$,
\[
\mathrm{Re}(\langle Ax, y \rangle_{\mathcal{H}}) = \frac{1}{4}\langle A(x + y), x + y \rangle_{\mathcal{H}} - \frac{1}{4}\langle A(x - y), x - y \rangle_{\mathcal{H}}
\]
成立. 因此由三角不等式与中线定理,
\begin{align*}
|\mathrm{Re}(\langle Ax, y \rangle_{\mathcal{H}})| &\leq \frac{1}{4}|\langle A(x + y), x + y \rangle_{\mathcal{H}}| + \frac{1}{4}|\langle A(x - y), x - y \rangle_{\mathcal{H}}| \\
&\leq \frac{1}{4}\sup_{u \in \mathcal{H}, \|u\| \leq 1}|\langle Au, u \rangle_{\mathcal{H}}|(\|x + y\|^{2} + \|x - y\|^{2}) \\
&\leq \sup_{u \in \mathcal{H}, \|u\| \leq 1}|\langle Au, u \rangle_{\mathcal{H}}|
\end{align*}
成立. 又, 另一方面, 由
\begin{align*}
&\sup_{x, y \in \mathcal{H}, \|x\| \leq 1, \|y\| \leq 1}|\mathrm{Re}(\langle Ax, y \rangle_{\mathcal{H}})| \\
&= \sup_{x, y \in \mathcal{H}, \theta \in \mathbb{R}, \|x\| \leq 1, \|y\| \leq 1}|\mathrm{Re}(\langle Ax, e^{i\theta}y \rangle_{\mathcal{H}})| \\
&= \sup_{x, y \in \mathcal{H}, \theta \in \mathbb{R}, \|x\| \leq 1, \|y\| \leq 1}|\mathrm{Re}(e^{-i\theta}\langle Ax, y \rangle_{\mathcal{H}})|
\end{align*}
得
\[
\sup_{x, y \in \mathcal{H}, \|x\| \leq 1, \|y\| \leq 1}|\mathrm{Re}(\langle Ax, y \rangle_{\mathcal{H}})| = \sup_{x, y \in \mathcal{H}, \|x\| \leq 1, \|y\| \leq 1}|\langle Ax, y \rangle_{\mathcal{H}}|
\]
成立. 因此由定理~\ref{thm:2.27}, 结论显然.
\end{proof}
